\documentclass[11pt,a4paper]{article}
\usepackage{amsmath,amssymb,geometry,booktabs,xcolor,setspace,enumitem,
            array,fancyhdr,titlesec,hyperref}
\geometry{margin=1in}
\onehalfspacing
\definecolor{accent}{RGB}{30,80,140}
\titleformat{\section}{\large\bfseries\color{accent}}{}{0em}{}[\titlerule]
\titleformat{\subsection}{\normalsize\bfseries}{}{0em}{}
\hypersetup{colorlinks,linkcolor=accent,urlcolor=accent}
\pagestyle{fancy}\fancyhf{}
\rhead{\small Kenya Time-Use Model}
\cfoot{\small\thepage}

\title{\textbf{Kenya Time-Use Model:\\ Spatial Equilibrium with Food Home Production}\\[6pt]
\large Model Specification, Calibration, and Counterfactuals}
\author{}\date{}

\begin{document}
\maketitle\thispagestyle{empty}
\tableofcontents\newpage

%──────────────────────────────────────────────────────────────────────────────
\section{Overview}
%──────────────────────────────────────────────────────────────────────────────

The model describes a household composed of a man and a woman in a Kenyan
county $\ell$. Each partner allocates time across market work, food
preparation, market care, market domestic services, home care, home
domestic services, and leisure. The household consumes food, non-food
goods, care services, and domestic services. Food and the two service
categories can be produced at home using the partners' time; non-food
goods are purchased only on the market.

This document describes the full extended model with the food/non-food
split. The previous specification (without the food split) treated food
as part of the residual goods category $x$.

%──────────────────────────────────────────────────────────────────────────────
\section{Household Problem}
%──────────────────────────────────────────────────────────────────────────────

\subsection{Time endowment}

Each partner $g \in \{m, f\}$ has total time normalised to 1 per period
(week). Time is allocated across:
\[
  L^M_g + L^{x_f}_g + L^c_g + L^d_g + \text{leisure}_g = 1
\]
where:
\begin{itemize}[noitemsep]
  \item $L^M_g$ : market work hours
  \item $L^{x_f}_g$ : food preparation hours \textbf{[new]}
  \item $L^c_g$ : home care hours
  \item $L^d_g$ : home domestic hours
\end{itemize}

\subsection{Consumption goods}

Four consumption categories:
\begin{itemize}[noitemsep]
  \item $x_f$ : food (home + market) \textbf{[new split]}
  \item $x_n$ : non-food goods (market only) \textbf{[new]}
  \item $c$ : care services (home + market)
  \item $d$ : domestic services (home + market)
\end{itemize}

\subsection{Home production}

For each home-producible good $i \in \{x_f, c, d\}$, household home
output is linear in time:
\[
  S^i_H = A^i_\ell \cdot \left( L^i_m + L^i_f \right)
\]
We set $A^i_\ell = 1$ for all $i$ and $\ell$ in the baseline, since
TFP is not separately identified from the data.

\subsection{Composite consumption}

For each home-producible category $i \in \{x_f, c, d\}$, the household
combines home and market via CES:
\[
  S^i = \left[ \omega_i (S^i_H)^{(\eta_i - 1)/\eta_i}
             + (1-\omega_i)(S^i_M)^{(\eta_i - 1)/\eta_i} \right]^{\eta_i / (\eta_i - 1)}
\]
Non-food $x_n$ is consumed directly without home production.

\subsection{Preferences (PIGL)}

Expenditure shares follow the Price-Independent Generalised Linearity
form. With four categories ($x_f, x_n, c, d$):
\[
  \vartheta^i = \beta_i + \kappa_i \left( \frac{E}{B} \right)^{-\varepsilon},
  \qquad \sum_i \beta_i = 1, \quad \sum_i \kappa_i = 0
\]
where $E$ is total expenditure and $B$ is the price index.

Here $\vartheta^{x_f}$ is the food share (typically $\approx 0.65$ in
Kenya, with $\kappa_{x_f} > 0$ since food is a necessity --- the share
falls with income), $\vartheta^{x_n}$ is non-food, etc.

\subsection{Disutility of work}

Each partner has a CES disutility aggregator over their four work
activities:
\[
  L^g = \left[
       D^g_M (L^M_g)^{(\rho-1)/\rho}
     + D^g_{x_f} (L^{x_f}_g)^{(\rho-1)/\rho}
     + D^g_c (L^c_g)^{(\rho-1)/\rho}
     + D^g_d (L^d_g)^{(\rho-1)/\rho}
     \right]^{\rho/(\rho-1)}
\]
with $\rho < 0$ (activities are complements in disutility).
The $D$ weights are gender-specific and capture both physical effort
and social norms.

\subsection{Indirect utility and participation}

Conditional on participation, the partner's value function is
\[
  V^g_P(h) = \phi \cdot w^g_\ell \cdot h + \text{const}
\]
where $h$ is the person's human capital and $\phi$ is the labour supply
elasticity scaling. Participation follows a logit rule:
\[
  P^g(h) = \frac{1}{1 + \exp\!\left[\,(\bar u^g - V^g_P(h))/\sigma^g_u\,\right]}
\]

%──────────────────────────────────────────────────────────────────────────────
\section{Spatial Equilibrium}
%──────────────────────────────────────────────────────────────────────────────

\subsection{County-level value function}

The representative county utility $V^*_\ell$ is the
population-weighted average of household value functions across the
human-capital distribution.

\subsection{Free-mobility condition}

In equilibrium households are indifferent across counties up to an
amenity term $\xi_\ell$:
\[
  V^*_\ell + \xi_\ell = \bar U \quad \forall \ell
\]
Counties with high wages must compensate with low amenities (negative
$\xi$); low-wage counties have high amenities. The amenity is identified
as the residual $\xi_\ell = \bar U - V^*_\ell$, calibrated so that
$\bar U$ equals the population-weighted mean of $V^*_\ell$.

%──────────────────────────────────────────────────────────────────────────────
\section{Calibration}
%──────────────────────────────────────────────────────────────────────────────

\subsection{Data sources}

\begin{itemize}[noitemsep]
  \item Kenya Time-Use Survey (TUS) 2021: time allocation by activity,
    24{,}004 individuals.
  \item Kenya Continuous Household Survey (KCHS): consumption aggregates
    (17{,}894 households), labour module (71{,}239 individuals),
    non-food items detail (279{,}356 rows).
\end{itemize}

\subsection{Engel curves with food/non-food split}

PIGL shares are estimated by WLS regression on the KCHS consumption
aggregate, using:
\begin{itemize}[noitemsep]
  \item $\vartheta^{x_f} = \texttt{padqfdcons} / E$ (food share, $\approx 0.65$)
  \item $\vartheta^{x_n} = (\texttt{padqnfitems} + \texttt{padqrent}
    + \texttt{padqeduc}) / E$ (non-food, $\approx 0.35$ minus services)
  \item $\vartheta^c, \vartheta^d$ : market care/domestic shares from
    the non-food items file (COICOP 13 and 5 respectively)
\end{itemize}

The slopes $\partial \vartheta^i / \partial \log E$ identify the
$\kappa_i$ coefficients. The food slope is negative (food is a
necessity); non-food, care, and domestic slopes are positive (luxuries).

\subsection{Time-use parameters (TUS)}

Food preparation hours $L^{x_f}$ come from the TUS code \texttt{t231}
(``Food and meals management and preparation''). The data shows striking
gender asymmetry:
\begin{itemize}[noitemsep]
  \item Men: 2.83h/week unconditional, 11.34h conditional (24\% participate)
  \item Women: 18.55h/week unconditional, 20.64h conditional (91\% participate)
\end{itemize}

Conditional ratio $r_{x_f} = L^{x_f}_m / L^{x_f}_f = 11.34 / 20.64 = 0.55$.

\subsection{Disutility weights via the solver fixed-point}

To avoid the sign error that previously plagued $D^f_M$, all
gender-specific disutility weights are derived from the solver's
fixed-point update map, not from a rearranged FOC.

For market work:
\[
  D^f_M = \frac{r_M}{\text{wage\_gap}^\rho}
  = \frac{1.474}{0.831^{-0.5}} = 1.344
\]

For food preparation, using the analogous home-activity update:
\[
  L^{x_f}_m / L^{x_f}_f = (D^f_{x_f} / D^m_{x_f})^\rho
  \quad\Rightarrow\quad
  D^f_{x_f} = D^m_{x_f} \cdot r_{x_f}^{1/\rho}
\]
With $D^m_{x_f} = 1$ (normalisation) and $r_{x_f} = 0.55$:
\[
  D^f_{x_f} = 1 \cdot 0.55^{-2} = 3.31
\]

This is large and negative--sign-corrected: women find food prep less
burdensome than men in the model not because they prefer it, but because
they spend so much time on it that the model needs $D^f_{x_f} > 1$ to
make the FOC consistent.

\subsection{Labour supply elasticity}

$\phi = 0.50$ is set as a literature prior. The IV estimate from KCHS
education-instrumented log wages gives $\hat\phi \approx 0.03$, but this
captures the uncompensated semi-elasticity (downward biased relative to
the Frisch elasticity needed in the model) and uses a weak instrument.
A range of 0.3--0.5 is standard in the macro labour literature; we use
0.5 as the upper end to keep the labour supply channel meaningful in
counterfactuals.

\subsection{Other parameters}

See \texttt{calibrated\_params.json} for the full list. Key ones:
\begin{itemize}[noitemsep]
  \item $\rho = -0.5$ (CES disutility curvature, literature prior)
  \item $\omega_i = 0.75, 0.70$ (CES home weights, priors)
  \item $\eta_i = 2.5$ (CES home/market substitutability, prior)
  \item $A^i_\ell = 1$ everywhere (home TFP not identified)
  \item $a = 0.2$ (1000-KSh non-labour income, hardcoded)
\end{itemize}

%──────────────────────────────────────────────────────────────────────────────
\section{Counterfactuals}
%──────────────────────────────────────────────────────────────────────────────

We study three counterfactuals:

\begin{enumerate}[leftmargin=*]
  \item \textbf{Closing the gender wage gap.} Set wage\_gap $= 1.0$
    in all counties; women now earn the same hourly wage as men.
  \item \textbf{Care price reduction.} Reduce $p^c_\ell$ by 30\% in
    all counties. Proxy for a national childcare subsidy or formalisation
    of the care sector.
  \item \textbf{Domestic price reduction.} Reduce $p^d_\ell$ by 30\%
    in all counties.
\end{enumerate}

For each counterfactual we compute, by county:
\begin{itemize}[noitemsep]
  \item $\Delta\text{GDP}_\ell$ (proxied by $N_\ell \cdot E_\ell$)
  \item $\Delta P^f_\ell$ (female participation)
  \item $\Delta(L^M_f / L^M_m)_\ell$ (gender hours ratio)
\end{itemize}

Results are presented as choropleth maps in \texttt{counterfactuals.pdf}.

%──────────────────────────────────────────────────────────────────────────────
\section{Known Limitations}
%──────────────────────────────────────────────────────────────────────────────

\begin{itemize}
  \item $N_\ell$ in \texttt{county\_fundamentals.csv} is the raw sum of
    TUS person-day weights ($\sim 10^8$), not the actual county population.
    This affects only absolute GDP levels; relative changes are unaffected.
  \item Three counties have spurious service prices from sample sparsity:
    Homa Bay ($p^d = 3$), Kakamega ($p^c = 351$), Tharaka-Nithi
    ($p^d = 110$). The model responds correctly to these prices
    (predicting more or less home production accordingly), but the
    prices themselves are unreliable.
  \item Home TFP $A^i_\ell$ is set to 1 everywhere because it cannot
    be separately identified from the data with current measurement.
  \item $\sigma_{\text{mig}} = 1.0$ in the spatial equilibrium is
    uncalibrated.
\end{itemize}

\end{document}
